\section{Introduction}

Quarkonium and leptonium systems constitute unique laboratories for investigating bound-state dynamics in quantum field theory.
Quarkonia probe both short- and long-distance effects of the strong interaction, making their production and decay mechanisms a long-standing testing ground for theoretical descriptions of hadronic bound states. The production and decay of quarkonia are studied in a variety of experimental environments, with measurements performed,
for instance, at electron-positron, lepton-hadron, and hadron-hadron colliders. For reviews, we refer the reader to refs.~\cite{Brambilla:2010cs,Lansberg:2019adr,Chapon:2020heu,Boer:2024ylx}. 
Leptonia, on the other hand, have not yet been observed at high-energy collider facilities. Positronium, the first predicted “onium” state in 1931~\cite{Mohorovicic1934}, was experimentally observed in 1951 through its formation in low-pressure gases~\cite{Deutsch:1951zza}.
Nowadays, these exotic systems serve, for instance, in low-energy experiments for precision studies of electromagnetic interactions~\cite{Karshenboim:2005iy} and for tests of the equivalence principle with antimatter~\cite{AEgIS:2023lpw}. Moreover, they can be exploited to investigate possible violations of discrete spacetime symmetries such as CPT invariance~\cite{Bernreuther:1988tt,Yamazaki:2009hp}. Here, C, P, and T denote charge conjugation, parity, and time reversal, respectively.

Non-relativistic QCD (NRQCD)~\cite{Bodwin:1994jh} is a rigorous low-energy effective field theory of QCD, which
enables systematic improvements in the calculation of quarkonium production and decay rates and remains, to date, the most widely adopted theoretical framework for quarkonium studies. Similarly, its QED counterpart, non-relativistic QED (NRQED)~\cite{Caswell:1985ui}, provides a framework for understanding leptonium states.
Within these frameworks, production cross sections are factorised into short-distance coefficients and long-distance matrix elements (LDMEs). Within these non-relativistic quantum field theories (NRQFTs), the spectroscopy, production rates, and decay rates of quarkonia and leptonia admit a controlled expansion in the relative velocity $v$ of their constituents.
This expansion organises the underlying interactions according to distinct power-counting schemes, commonly referred to as velocity-scaling rules.
While the so-called \swavetxt\ configurations, characterised by vanishing intrinsic orbital angular momentum, are among the most widely studied states both theoretically and experimentally, they are not the only ones of interest. States with non-vanishing orbital angular momentum provide complementary probes to test the factorisation of the employed NRQFTs and the validity of the velocity expansion. Among these, we focus in this work on \pwavetxt\ states, which possess one unit of orbital angular momentum.

In heavy-quarkonium physics, \pwavetxt\ states offer direct sensitivity to spin-orbit and tensor interactions in the heavy-quark potential and have long served as precision probes of quarkonium spectroscopy. Beyond this, they also play a central role in the theoretical description of quarkonium production within NRQCD. 
For instance, the transverse-momentum spectrum of inclusive $J/\psi$ hadroproduction receives significant contributions from colour-octet (CO) \pwavetxt\ channels, despite their LDMEs being suppressed relative to the colour-singlet (CS) \swavetxt\ state.
Moreover, the production of \pwavetxt\ states highlights important subtleties of NRQCD factorisation, such as the cancellation of infrared (IR) divergences between short-distance coefficients and LDMEs~\cite{Bodwin:1992ye,Bodwin:1994jh,Petrelli:1997ge,AH:2024ueu}.
Experimentally, \swavetxt\ quarkonium yields may receive significant feed-down contributions from decays of \pwavetxt\ states. A precise description of these yields therefore requires accurate modelling of \pwavetxt\ production and decay mechanisms.
{\pwavetxt}s are thus indispensable for understanding inclusive quarkonium production in high-energy particle collisions and remain a key ingredient in global fits of CO LDMEs~\cite{Ma:2010vd,Ma:2010yw,Butenschoen:2010rq,Ma:2010jj,Butenschoen:2011yh,Butenschoen:2012px,Chao:2012iv,Wang:2012is,Gong:2012ug,Shao:2012fs,Butenschoen:2012qr,Gong:2013qka,Shao:2014fca,Bodwin:2014gia,Shao:2014yta,Han:2014kxa,Feng:2015wka,Bodwin:2015iua,Feng:2018ukp,Feng:2020cvm,Brambilla:2021abf,Brambilla:2022rjd,Butenschoen:2022qka,Brambilla:2022ayc,Brambilla:2024iqg}.

In the QED sector, leptonium systems, such as positronium, muonium, and their heavier siblings, provide a complementary and theoretically cleaner arena in which analogous \pwavetxt\ states can be studied. Governed solely by QED, these systems allow for high-precision theoretical predictions without the complications of confinement or non-perturbative hadronic effects. Experimentally, \pwavetxt\ leptonic bound states have already been observed in low-energy experiments~\cite{PhysRevLett.34.177,PhysRevLett.102.133202}. Theoretically, their properties are well understood within NRQED and can be calculated with high precision~\cite{dEnterria:2022alo}, making them promising probes of higher-order QED radiative corrections, fundamental Standard Model (SM) constants such as the fine-structure constant and lepton masses~\cite{dEnterria:2023yao}, as well as potential physics beyond the SM (BSM).

Given their relevance, public tools that support \pwavetxt\ generation to some extent are available. However, owing to the increased theoretical and computational complexity of \pwavetxt\ calculations, such tools are fewer in number than those publicly available for \swavetxt\ production. In general-purpose Monte Carlo (MC) event generators, \pwavetxt\ quarkonia can be generated using parton showers in \pythia~{\tt v8.3}~\cite{Cooke:2023ldt} and \herwig~{\tt v7.4}~\cite{Masouminia:2025kec}. In addition, \pythia\ provides matrix elements for a limited set of $2\to 2$ single \pwavetxt\ quarkonium production processes~\cite{Sjostrand:2000wi,Sjostrand:2006za,Sjostrand:2007gs,Sjostrand:2014zea,Bierlich:2022pfr}. 
More specialised implementations include exclusive double-diffractive $\chi_{Q}$ production in the \texttt{SuperChic} generator~\cite{Khoze:2000jm,Harland-Lang:2009kjv}, as well as simulations based on the phenomenological Wigner density matrix formalism in \texttt{EPOS4}~\cite{Song:2017phm,Villar:2022sbv,Zhao:2025cnp}, which do not rely on NRQCD factorisation. A dedicated parton-level generator, \texttt{BCVEGPY}~{\tt v2.0}~\cite{Chang:2005hq}, provides leading-order (LO) simulations of inclusive hadroproduction of excited \pwavetxt\ $B_c$ states, while \texttt{FDCHQHP}~\cite{Wan:2014vka} enables next-to-leading order (NLO) calculations of differential cross sections for inclusive \swave- and \pwavetxt\ charmonium and bottomonium production at hadron colliders.
Besides these largely process-specific tools, two process-independent implementations are available, \madonia~\cite{Artoisenet:2007qm} and \helaconia~\cite{Shao:2012iz,Shao:2015vga}. 
They provide parton-level event generation for processes involving \swave- and \pwavetxt\ quarkonia at tree level within NRQCD. 
\madonia, based on \texttt{MadGraph/MadEvent}~{\tt v4}~\cite{Alwall:2007st} and no longer actively maintained, is restricted to final states containing at most one \swave- or \pwavetxt\ quarkonium. By contrast, \helaconia, built upon \texttt{HELAC-PHEGAS}~\cite{Kanaki:2000ey,Papadopoulos:2000tt,Kanaki:2000ms,Papadopoulos:2006mh}, is more flexible and imposes no such restriction on the number of final-state quarkonia.\footnote{In the current public versions of \helaconia, the number of \pwavetxt\ quarkonia cannot exceed two, while no restriction applies to \swavetxt\ states. This limitation is expected to be lifted in future versions.}
Both programs generate unweighted LO events that can be interfaced with general-purpose MC generators via the Les Houches Event (LHE) file format~\cite{Alwall:2006yp}. Finally, there are also public software packages, such as {\tt FeynOnium}~\cite{Brambilla:2020fla} and {\tt AmpRed}~\cite{Chen:2024xwt}, for semi-automatic symbolic or analytical calculations of matrix elements  within NRQFTs.

Despite these considerable efforts, the aforementioned tools still present important limitations, the most common of which is their restriction to LO calculations. This is a notable shortcoming, since LO predictions are often insufficient to meet the precision requirements of modern experiments, and higher-order perturbative corrections can be sizeable. While process-by-process calculations remain possible, they lack the high degree of automation achieved for processes involving only elementary particles. Automated tools are therefore highly desirable to simplify calculations involving bound states.
Within the \madgraph\ framework~\cite{Alwall:2014hca,Frederix:2018nkq} (hereafter referred to as \mgshort), a programme is underway to overcome such limitations, aiming to develop an all-encompassing tool for bound-state production studies. The first necessary step is to extend \mgshort\ to simulate arbitrary \swave- and \pwavetxt\ quarkonium and leptonium production processes at tree-level within NRQCD and NRQED.
The programme was initiated in an earlier work~\cite{ColpaniSerri:2025vdz}, where the \swavetxt\ implementation was presented. 
Here, we upgrade this framework by automating the generation of events with an arbitrary number of \pwavetxt\ states, thereby completing the LO automation of \swave- and \pwavetxt\ quarkonium and leptonium production processes.
Our implementation of the \pwavetxt\ states and the associated projectors relies on dual numbers~\cite{Clifford:1871aaa}; this choice differs from those adopted in the literature and is, a priori, well suited for further extensions towards NLO calculations.
The \swave- and \pwavetxt\ implementations are combined within a unified automated framework, denoted as the \textcolor{teal}{\madsons} (\texttt{\textcolor{teal}{Mad}Graph} \textcolor{teal}{S}imulations \textcolor{teal}{O}f \textcolor{teal}{N}RQFT \textcolor{teal}{S}tates) module.

As a second novel development in this module, we introduce an improved treatment of quarkonium masses, which are conventionally approximated within NRQCD by the sum of the masses of the heavy quark-antiquark pair. 
This approximation is required to preserve fundamental properties of the $S$-matrix in the calculation of short-distance amplitudes.
To account for the difference between the physical quarkonium mass and the sum of the constituent masses, we propose to retain the physical quarkonium masses during phase-space generation. The resulting phase space is then mapped onto parton-level four-momenta used for matrix-element evaluation through a momentum-reshuffling procedure, a technique widely used in high-energy MC tools, for instance in parton-shower algorithms~\cite{Bengtsson:1987kr,Marchesini:1987cf}. This strategy offers several advantages. 
First, strictly speaking, a consistent treatment of physical mass effects would require the inclusion of higher-order relativistic corrections in both the matrix element and the phase space. While their inclusion in the former is highly non-trivial, the present method provides a simple way to incorporate them in the latter, thereby improving MC simulations near threshold with minimal effort. Second, without momentum reshuffling, generated events assign a common mass to all quarkonia composed of the same heavy quarks.\footnote{For instance, in $J/\psi + \psi(2\swave)$ production, the two physical masses differ by approximately $590\,{\rm MeV}$, yet both quarkonium masses are taken to be twice the chosen charm-quark mass.} As a result, variations in the constituent-quark mass can introduce sizeable uncertainties in the predicted cross sections. Fixing the phase space through momentum reshuffling is, a priori, expected to reduce this dependence, particularly in threshold regions. Finally, the approach automatically accounts for mass differences between the produced quarkonia, feed-down contributions from heavier states, and CO intermediate states, which can be taken heavier than the corresponding quarkonium state, thereby matching, for instance, the hadronisation models used in \pythia{\tt 8}. This, in turn, enhances the automation of interfaces to general-purpose MC generators.

The remainder of the paper is organised as follows. In section~\ref{sec:nrqft}, we review the NRQFT factorisation framework underlying NRQCD and NRQED, which provides the foundation for our implementation. The technical details of the implementation are presented in section~\ref{sec:implementation}. Validation tests and benchmark studies are discussed in section~\ref{sec:benchmark}, followed by phenomenological applications to quarkonium and leptonium production in sections~\ref{sec:quarkonium} and~\ref{sec:leptionium}, respectively. Our conclusions are given in section~\ref{sec:conclusions}. Unless explicitly stated otherwise, all numerical computations presented in this work are carried out using the default setup described in appendix~\ref{app:setup}.